\uuid{pkaV}
\titre{Ensembles}
\auteur{bodin ; croizat ; sacre}
\organisation{exo7}
\licence{CC BY-NC-SA 4.0}
\datecreate{2022-03-08}
\niveau{L1}
\module{Algèbre}
\chapitre{Logique, ensemble, raisonnement}
\sousChapitre{Ensemble}
\difficulte{2}
\theme{Ensembles}
\exo7qcmid{hilisit.logique.032}

\contenu{
\begin{question}
\qclassification{100.02}
Soit l'ensemble $A = \{ -1 , 0 , 1 \}$ et la fonction réelle donnée par $f(x) = x^2-1$. Quelles sont les assertions vraies ?
\begin{answers}
	\bad{$\forall x \in A \quad f(x) = 0$}
	\good{$\exists \, x \in A \quad f(x) = 0$}
	\bad{$f$ est bijective de $A$ dans son image $f(A)$.}
	\good{$\forall x \in A \quad f(x) \in A$}
\end{answers}
\begin{explanations}
On a $f(-1) = f(1) = 0 \in A$, $f(0) = -1 \in A$ mais $-1 \neq 0$.\\
La fonction $f$ n'est pas bijective de $A$ dans $f(A)$ puisque $1$ et $-1$ ont la même image.
\end{explanations}
\end{question}
}
